% Source: report/sections/07_generalisation.md (owner: Manoj, #6).
\section{Generalisation to CVC-ClinicDB}
\label{sec:gen}
A model that will be used in another hospital must work on images from other cameras, other lighting and other
patients. To test this, every model trained on Kvasir-SEG was also evaluated on \textbf{CVC-ClinicDB}~\cite{bernal2015cvc}
\textbf{without any retraining}: 612 frames from 29 colonoscopy videos recorded in Barcelona with a different
endoscope (\cref{sec:extra-data}). The images are resized and normalised exactly like the Kvasir-SEG images, and the
threshold is 0.5.

\begin{table}[t]
  \centering
  \small
  \setlength{\tabcolsep}{3pt}
  \begin{adjustbox}{max width=\linewidth}
  \begin{tabular}{@{}lcccc@{}}
    \toprule
    Model & Kvasir Dice ($n{=}120$) & CVC Dice ($n{=}612$) & CVC IoU & Drop \\
    \midrule
    U-Net (scratch)     & 0.851 & 0.750 & 0.655 & 0.101 \\
    U-Net + ResNet-34   & 0.903 & 0.828 & 0.754 & 0.075 \\
    U-Net + MobileNetV2 & 0.894 & 0.753 & 0.680 & 0.141 \\
    Attention U-Net     & 0.854 & 0.741 & 0.646 & 0.113 \\
    SegFormer (MiT-B2)  & 0.917 & 0.863 & 0.787 & 0.054 \\
    \bottomrule
  \end{tabular}
  \end{adjustbox}
  \caption{Dice on the Kvasir-SEG test set and on CVC-ClinicDB, and the drop between them.}
  \label{tab:cvc}
\end{table}

\paragraph{Models trained from scratch.} The baseline loses 0.101 Dice on the new hospital's images ($0.851 \rightarrow 0.750$, \cref{tab:cvc}). It
still finds most polyps, but a tenth of its overlap is lost, which shows that it has partly learned features
specific to Kvasir-SEG's camera and image style. The Attention U-Net loses slightly more (0.113) and ends up a
little below the baseline on CVC (0.741 against 0.750): the attention gates do not make the model more robust
to a new hospital either.

\paragraph{SegFormer.} SegFormer loses only 0.054 Dice, about half as much as the two models trained from scratch, and
keeps a Dice of 0.863 on CVC-ClinicDB. Its lead over the baseline grows from 0.066 on Kvasir-SEG to 0.113 on
CVC-ClinicDB. Two of its properties may explain this: its encoder was pretrained on ImageNet, so it starts from
features learned on over a million varied photographs rather than on 800 colonoscopy images; and its
self-attention relates every part of the image to every other part, which may make it rely less on local colour
and texture that change between cameras.

\paragraph{Pretrained CNN encoders.} ResNet-34 and MobileNetV2 were also pretrained on ImageNet but are CNNs, so they help
separate these two explanations. They do not behave alike. ResNet-34 loses 0.075 Dice and keeps 0.828 on
CVC-ClinicDB, between the scratch models and SegFormer. MobileNetV2 loses 0.141, \textbf{the largest drop of all five
models}: on Kvasir-SEG it was 0.043 above the baseline, on CVC-ClinicDB only 0.003. ImageNet pretraining alone
therefore does not guarantee robustness to a new hospital. It helped the larger ResNet-34 encoder (24.4\,M
parameters), but not the compact MobileNetV2 (6.6\,M), whose small capacity may be used up fitting Kvasir-SEG's
own image style. SegFormer, pretrained \emph{and} attention-based, transfers best, which suggests that both pretraining
and the transformer encoder contribute, although one seed per model and one external dataset cannot prove this.

The ranking on CVC-ClinicDB is almost the same as on Kvasir-SEG: SegFormer, then ResNet-34, then the other three.
The one change is that MobileNetV2 falls back to the level of the scratch models (0.753 against 0.750 and 0.741).
For deployment in a hospital whose images were not seen in training, SegFormer and ResNet-34 are the safer
choices.

\paragraph{Caveats.}
\begin{itemize}
  \item CVC-ClinicDB frames come from only 29 videos, so neighbouring frames are very similar and the 612 frames are
        far from 612 independent cases.
  \item Its images are smaller ($384\times288$) than most Kvasir-SEG images and wider than they are tall. Resizing them to
        $352\times352$ stretches them vertically and changes their apparent sharpness, so polyps look slightly different
        from those the models were trained on.
  \item We use no CVC image for training or for choosing anything, so these numbers are a fair test of transfer, but a
        single seed means differences of about 0.01 (such as between the baseline, the Attention U-Net and MobileNetV2
        on CVC) should not be over-interpreted.
\end{itemize}
